\documentclass[12pt]{article}
\usepackage[utf8]{inputenc}
\usepackage[T1]{fontenc}
\usepackage{hyperref}
\title{Istro Romanian Language Resources for K-12 Curriculum Development}
\author{Compiled by Claude Sonnet 5.5 (claude-sonnet-5-5) for Andr\'as Kornai}
\date{}
\begin{document}
\maketitle

\section{Introduction}
Istro-Romanian (ISO 639-3 ruo) is a severely endangered Balkan Romance language spoken in six mountain villages in Istria, Croatia. Speakers in \v{S}u\v{s}njevica call it Vlashki, and speakers in \v{Z}ejane call it Zheyanski. Search found a dictionary, a documentation project and some history, but no current school curriculum.

\section{Primary Educational Resources}
A school called ``Traian Emperor'' ran in \v{S}u\v{s}njevica from 1921 to 1925. It was founded by Andrei Glavina and taught in both the Istro-Romanian dialect and literary Romanian. This is a historical precedent for local-language schooling: \url{https://en.wikipedia.org/wiki/Andrei_Glavina}

\section{Linguistic and Reference Materials}
\begin{itemize}
\item Kova\v{c}ec's 1998 Istro-Romanian-Croatian dictionary, online edition. It was the only published bilingual dictionary at the time of the search result: \url{https://www.vlaski-zejanski.com/en/o-rjecniku}
\item hISTROX, Oxford: Istro-Romanian language, \v{Z}ejanski and Vlaški: \url{https://istrox.ling-phil.ox.ac.uk/istro-romanian-language-zejanski-and-vlaski}
\item hISTROX news and events, which mentions the fifth volume of the Romanian Academy's Dictionary of the Istro-Romanian Dialect: \url{https://histrox.ling-phil.ox.ac.uk/article/news-and-events}
\item University of Rijeka documentation project (Z. Vrzi\'c): a trilingual Vlashki/Zheyanski-Croatian-English dictionary, annotated corpus and syntax description are planned: \url{https://cji.uniri.hr/en/?p=543}
\end{itemize}

\section{Religious and Community Resources}
Community website on Istro-Romanian: \url{https://istro-romanian.net/articles/art090708.html}

\section{Digital and Audio Resources}
The online dictionary (above) is the main digital resource found. Whether the Rijeka project has published audio or corpus material was not checked.

\section{Curriculum Development Context}
Only adults speak the language as a first language, and it is not known to be taught in schools. A curriculum would be a revitalization programme built on the Kova\v{c}ec dictionary and the Rijeka and Oxford documentation work, with speaker involvement.
\end{document}